
\assignq{2}{2}{Superdense coding and a passive linear-optical Bell measurement}{2--3}
Start with \(\ket{\Phi^+}_{AB}\) and act on
\emph{Alice's} qubit only. The problem defines
\(X\ket0=\ket1,\ X\ket1=\ket0\),
\(Z\ket0=\ket0,\ Z\ket1=-\ket1\).

\subsection*{(a) Calculate the four encoded Bell states}
Direct application to each computational term gives
\[
\begin{array}{c|c|c}
\text{message}&\text{gate on }A&\text{output state}\\\hline
00&I&(\ket{00}+\ket{11})/\sqrt2=\ket{\Phi^+}\\
01&X&(\ket{10}+\ket{01})/\sqrt2=\ket{\Psi^+}\\
10&Z&(\ket{00}-\ket{11})/\sqrt2=\ket{\Phi^-}\\
11&XZ&(\ket{10}-\ket{01})/\sqrt2=-\ket{\Psi^-}
\end{array}
\]
The last global minus phase is unobservable, so
the four encodings specify four orthogonal Bell states.

\subsection*{(b) Why two bits can be sent}
Before communication, Alice and Bob share
one entangled Bell pair. Alice chooses one
of four unitaries (\(2^2\) alternatives), sends
her \emph{one} photon, and Bob performs an
ideal joint Bell-basis measurement.
Four orthogonal outputs encode two classical bits.
The shared entanglement and the transmitted
physical qubit together provide the resource;
this is not two bits from an arbitrary isolated
one-photon qubit.

\subsection*{(c) Expand the specified beam-splitter measurement}
For each polarisation label \(\lambda=H,V\),
\[
a_\lambda^\dagger\longrightarrow
\frac{c_\lambda^\dagger+d_\lambda^\dagger}{\sqrt2},
\qquad
b_\lambda^\dagger\longrightarrow
\frac{c_\lambda^\dagger-d_\lambda^\dagger}{\sqrt2}.
\]
For two equal polarisations,
\(a_\lambda^\dagger b_\lambda^\dagger
\to[(c_\lambda^\dagger)^2-
(d_\lambda^\dagger)^2]/2\). Therefore
\[
\boxed{\ket{\Phi^\pm}\longrightarrow
\frac{(c_H^\dagger)^2-(d_H^\dagger)^2
\pm(c_V^\dagger)^2\mp(d_V^\dagger)^2}
{2\sqrt2}\ket{\mathrm{vac}}}.
\]
For the \(\Psi\) states the mixed-polarisation
amplitudes interfere:
\[
\boxed{\ket{\Psi^+}\longrightarrow
\frac{c_H^\dagger c_V^\dagger
-d_H^\dagger d_V^\dagger}{\sqrt2}
\ket{\mathrm{vac}}},\qquad
\boxed{\ket{\Psi^-}\longrightarrow
\frac{c_V^\dagger d_H^\dagger
-c_H^\dagger d_V^\dagger}{\sqrt2}
\ket{\mathrm{vac}}}.
\]
A polarising beam splitter and number-resolving detectors
at each output distinguish these patterns:
\[
\begin{array}{c|l|c}
\text{Bell input}&\text{Detection signature}&
\text{Identification}\\\hline
\Phi^+&2\text{ photons same port, same pol.}&
\text{ambiguous}\\
\Phi^-&2\text{ photons same port, same pol.}&
\text{ambiguous}\\
\Psi^+&2\text{ orthogonal pol., same port}&
\text{unique}\\
\Psi^-&1\text{ photon each port, orthogonal pol.}&
\text{unique}
\end{array}
\]
The relative signs between \(\Phi^\pm\) components
do not change these photon-count patterns.

\subsection*{(d) Decoding efficiency}
With four equiprobable messages, only
\(\Psi^+,\Psi^-\) are unambiguous:
\[
\boxed{P_{\mathrm{unambiguous}}=\frac24=\frac12}.
\]
Thus the standard passive linear-optical analyser
alone cannot implement a \emph{deterministic}
four-outcome Bell measurement. Extra
ancillary resources or different physical
interactions may improve the measurement.
\methodnote{Bosonic operators for distinct modes commute.
Add coherent creation-operator amplitudes before converting
into Fock-state probabilities.}
